\documentclass[11pt,a4paper]{article}
\usepackage[T1]{fontenc}
\usepackage[utf8]{inputenc}
\usepackage{geometry}
\usepackage{hyperref}
\usepackage{longtable}
\usepackage{booktabs}
\usepackage{xcolor}
\usepackage{catchfile}
\geometry{margin=25mm}
\hypersetup{
  colorlinks=true,
  linkcolor=blue,
  urlcolor=blue,
  pdftitle={FMDFlashcard Release Documentation},
  pdfauthor={FMDFlashcard contributors}
}

\CatchFileDef{\FMDVersion}{../../VERSION}{\endlinechar=-1 }
\newcommand{\command}[1]{\texttt{#1}}

\title{FMDFlashcard\\Release Documentation}
\author{FMDFlashcard contributors}
\date{Version \FMDVersion}

\begin{document}
\maketitle
\tableofcontents

\section{Purpose}

FMDFlashcard is a local-first desktop application for Markdown-based flashcards,
exams, and spaced-repetition study. The source documentation is maintained in
the repository's \command{docs/} directory and built strictly with MkDocs. This
PDF is release evidence: it records the supported package contract, integrity
checks, and maintainer-facing installation guidance for the corresponding
release.

\section{Supported desktop packages}

\begin{longtable}{@{}lll@{}}
\toprule
Platform & Architecture & Packages \\
\midrule
Windows & x86\_64 & MSI, NSIS setup executable, portable ZIP \\
Linux & x86\_64 & Debian package, RPM package, AppImage \\
macOS & Apple Silicon & DMG, compressed application bundle \\
macOS & Intel & DMG, compressed application bundle \\
\bottomrule
\end{longtable}

Every production package is built on a native GitHub-hosted runner. The
Linux-to-Windows cross-build command is an explicitly experimental local
fallback and is not accepted as a production Windows release source.

\section{Install and run from source}

From the repository root, inspect the host and install project dependencies:

\begin{verbatim}
./control doctor
./control install --dry-run
./control install
./control run --foreground
\end{verbatim}

PowerShell users can use \command{.\textbackslash control.ps1}; Command Prompt
users can use \command{control.cmd}. The canonical implementation remains
\command{python tools/control.py}.

\section{Quality and tests}

The complete local gates are exposed through stable commands:

\begin{verbatim}
./control quality
./control test --suite all --report --ci
./control docs check
./control version check
./control release check
\end{verbatim}

The frontend, Rust crate, Python tooling, documentation, workflow contracts,
and platform-independent Tauri plans are checked independently so a failure
cannot be hidden by another suite.

\section{Artifact integrity}

A release contains \command{release-manifest.json}, \command{SHA256SUMS}, and
\command{SBOM.spdx.json}. Verify downloaded files from the asset directory:

\begin{verbatim}
sha256sum --check SHA256SUMS
python tools/control.py release verify --directory .
\end{verbatim}

The manifest records each artifact's file name, size, SHA-256 digest, build
target, architecture, runner family, commit, signature state, notarization
state, and toolchain versions. GitHub provenance attestations can additionally
be checked with the GitHub CLI:

\begin{verbatim}
gh attestation verify <asset> --repo kleiveist/FMDFlashcard
\end{verbatim}

\section{Signing status}

Unsigned CI packages remain useful for pull-request and smoke validation.
Published assets state whether they are signed and, for macOS, whether
notarization was verified. A release configured with the required signing
policy fails before publication when credentials or verification evidence are
missing. Updater signing is intentionally absent because the application does
not implement the Tauri updater.

\section{Further documentation}

The full user and developer documentation is available in the release source
and at the repository:
\begin{itemize}
  \item \url{https://github.com/kleiveist/FMDFlashcard/tree/main/docs/usr}
  \item \url{https://github.com/kleiveist/FMDFlashcard/tree/main/docs/dev}
  \item \url{https://github.com/kleiveist/FMDFlashcard/tree/main/docs/tools}
\end{itemize}

\end{document}
